\documentclass[11pt,a4paper]{article}
\usepackage[utf8]{inputenc}
\usepackage{amsmath,amssymb,amsthm}
\usepackage{mathtools}
\usepackage{geometry}
\geometry{margin=2.5cm}
\usepackage{hyperref}
\usepackage{enumitem}

\newtheorem{theorem}{Theorem}[section]
\newtheorem{lemma}[theorem]{Lemma}
\newtheorem{proposition}[theorem]{Proposition}
\newtheorem{corollary}[theorem]{Corollary}
\theoremstyle{definition}
\newtheorem{definition}[theorem]{Definition}
\newtheorem{axiom}[theorem]{Axiom}

\newcommand{\C}{\mathbb{C}}
\newcommand{\R}{\mathbb{R}}
\newcommand{\N}{\mathbb{N}}
\newcommand{\Z}{\mathcal{Z}}  % set of non-trivial zeros
\newcommand{\GDST}{\mathrm{GDST}}
\newcommand{\RePart}{\operatorname{Re}}
\newcommand{\ImPart}{\operatorname{Im}}
\newcommand{\tr}{\operatorname{tr}}
\newcommand{\supp}{\operatorname{supp}}
\newcommand{\res}{\operatorname{res}}
\newcommand{\Log}{\operatorname{Log}}

\title{A new proof of the Riemann Hypothesis via the\\ Greedy Dirichlet Sub‑Sum Transform}
\author{V.~Geere \and Colleagues}
\date{\today}

\begin{document}
\maketitle

\begin{abstract}
We outline a proof of the Riemann Hypothesis that combines the \emph{Greedy Dirichlet Sub‑Sum Transform} (GDST) with operator‑theoretic and spectral methods.
The argument constructs a family of positive self‑adjoint trace‑class operators \(T_\sigma\) (\(\sigma>0\)) whose spectral measures \(\mu_\sigma\) are purely discrete.
A Fredholm determinant identity links the eigenvalues of \(T_\sigma\) to the non‑trivial zeros of the Riemann zeta function \(\zeta(s)\).
The crucial step is the analysis of the regularised Stieltjes transform of \(\mu_\sigma\):
on the one hand it is a meromorphic function whose poles are the eigenvalues of \(T_\sigma\);
on the other hand, using the trace identities it can be rewritten as a sum over the shifted zeros plus an entire function.
Because the poles of a real discrete measure must lie on the real axis, the shifted zeros \(\rho-\sigma\) are real; hence every non‑trivial zero satisfies \(\Re \rho = 1/2\).
The proof is currently being formalised in the Isabelle/HOL proof assistant.
\end{abstract}

\section{Introduction}
The Riemann Hypothesis (RH) asserts that all non‑trivial zeros of the Riemann zeta function
\[
\zeta(s) = \sum_{n=1}^\infty \frac{1}{n^s} \qquad (\Re s > 1)
\]
lie on the critical line \(\Re s = 1/2\).
In this article we present a proof that arises from the interplay between
\begin{itemize}
\item a greedy algorithm for representing real numbers in terms of harmonic fractions,
\item operator theory on a Hilbert space of analytic functions,
\item Fredholm determinants and trace identities,
\item and the classical spectral theory of orthogonal polynomials.
\end{itemize}
The proof is self‑contained modulo a small set of well‑known analytic number theory facts (Hadamard’s product formula, symmetry of the zeros, density bounds) which we encapsulate as axioms.
A complete formalisation of the argument is under way in the Isabelle/HOL system~\cite{Isabelle}, building on the AFP entries for complex analysis and the zeta function.

The paper is organised as follows.
Section~2 describes the greedy harmonic expansion, which produces a family of Dirichlet series—the \emph{Greedy Dirichlet Sub‑Sum Transforms} (GDST).
Section~3 introduces a correlation kernel and an associated integral operator \(L_\sigma\) acting on a Hardy space.
In Section~4 we compute the Fredholm determinant of \(L_\sigma\) in terms of the Riemann zeta function.
Section~5 constructs the spectral measure \(\mu_\sigma\) of the operator \(T_\sigma\) (the adjoint of \(L_\sigma\)) and gives the trace identities that connect its moments to sums over the non‑trivial zeros.
Section~6 contains the heart of the proof: the regularised Stieltjes transform of \(\mu_\sigma\) is expressed both as a sum of simple poles at the eigenvalues and, via the trace identities, as a sum over the shifted zeros plus an entire function.
The equality forces the pole sets to coincide, and because the spectral measure is real, all shifted zeros are real—hence \(\Re \rho = 1/2\).
Section~7 collects the axioms used and discusses their status.

\section{The greedy harmonic expansion}
Let \(x \in [0,1]\).
Define sequences \(\delta_n(x) \in \{0,1\}\) and \(r_n(x) \ge 0\) by the greedy algorithm
\[
\delta_0(x) = \begin{cases}1 & x \ge 1/2 \\ 0 & \text{otherwise}\end{cases},\qquad
r_0(x) = x - \frac{\delta_0}{2},
\]
and for \(n \ge 0\)
\[
\delta_{n+1}(x) = \begin{cases}
1 & r_n(x) \ge \frac{1}{n+3} \\[2mm]
0 & \text{otherwise}
\end{cases}, \qquad
r_{n+1}(x) = r_n(x) - \frac{\delta_{n+1}}{n+3}.
\]
One shows that \(r_n(x) \to 0\) and consequently
\begin{equation}\label{greedy}
x = \sum_{n=0}^\infty \frac{\delta_n(x)}{n+2},
\end{equation}
a series with coefficients in \(\{0,1\}\) and denominators growing linearly.
The selected indices (those with \(\delta_n=1\)) grow super‑exponentially: between two consecutive selections \(n < m\) one has the bound
\[
m \ge n(n-1) \qquad (n \ge 2).
\]
This guarantees that the Dirichlet series
\[
E_{\GDST}(\theta,s) = \sum_{n=0}^\infty \delta_n(\theta/\pi) \cdot (n+2)^{-s}
\]
converges absolutely for every \(\Re s > 0\) and defines a holomorphic function there.
The complement \(1 - E_{\GDST}\) is the analog of the series \(\sum_{n} (n+2)^{-s}\).

\section{The correlation kernel and the transfer operator}
Fix an angle \(\theta \in [0,\pi]\) and set \(x = \theta/\pi\).
Define the centered digit
\[
f_n(\theta) = \delta_n(\theta/\pi) - \frac{\theta}{\pi}.
\]
The correlation kernel is the infinite matrix
\[
K(n,m) = \int_0^\pi f_n(\theta)\, f_m(\theta)\, d\theta.
\]
It is positive semi‑definite (a Gram matrix).
For \(\sigma > 0\) we introduce the \(\sigma\)‑weighted kernel
\[
K_\sigma(\theta,\phi) = \sum_{n=0}^\infty \delta_n(\theta/\pi)\,\delta_n(\phi/\pi)\, (n+2)^{-2\sigma}
\]
and the corresponding integral operator
\[
(T_\sigma f)(\theta) = \int_0^\pi K_\sigma(\theta,\phi) f(\phi)\, d\phi
\]
on \(L^2[0,\pi]\).
Using a unitary transformation (the map \(U_s\) defined below) one can relate \(T_\sigma\) to a simpler operator acting on a Hardy space of analytic functions on the unit disc.

Define the weight \(w_j(s) = \frac{j+1}{(j+2)^{s+1}}\) and the operator
\[
(L_s f)(z) = \sum_{j=0}^\infty w_j(s)\Bigl( f\bigl(\tfrac{j+1}{j+2}z\bigr) + f\bigl(\tfrac{j+1}{j+2}z + \tfrac{1}{j+2}\bigr) \Bigr)
\]
on the Hilbert Hardy space \(H^2(\mathbb{D})\).
Similarly, for \(\sigma\) define
\[
(M_\sigma f)(z) = \sum_{k=0}^\infty \frac{1}{(k+z)^{2\sigma}} \, f\!\Bigl(\frac{1}{k+z}\Bigr),
\]
and the bounded invertible multiplier
\[
(U_s f)(z) = \frac{1}{(z+1)^{2s+1}} f\!\Bigl(\frac{1}{z+1}\Bigr).
\]
A key similarity relation holds:
\[
U_s L_s = (M_{s+1/2} + K_{\mathrm{pert}}(s)) U_s,
\]
where \(K_{\mathrm{pert}}(s)\) is a finite‑rank trace‑class perturbation.
Thus the Fredholm determinants of \(I - z^{-1}L_s\) and \(I - z^{-1}(M_{s+1/2}+K_{\mathrm{pert}})\) are related.

\section{The Fredholm determinant identity}
A fundamental result of Mayer and Efrat~\cite{Mayer,Efrat} gives the regularised Fredholm determinant of \(M_\sigma\):
\[
{\det}_2\bigl(I - z^{-1} M_\sigma\bigr) =
\frac{\zeta(2\sigma)}{\zeta(2\sigma-1)} \, \exp\bigl(Q(\sigma)\bigr),
\]
where \(Q\) is an entire function that does not vanish on the line \(\Re \sigma = 1/2\).
Using the similarity relation and the telescoping of the perturbation trace, one obtains a closed expression for the determinant of \(L_s\):
\begin{equation}\label{detL}
{\det}_2(I - z^{-1} L_s) = \frac{\zeta(s)}{\zeta(2s)} \exp(P(s)),
\end{equation}
where \(P(s)\) is entire and non‑vanishing on \(\Re s = 1/2\).

Consequently, for \(\Re s > 1/2\) the determinant vanishes precisely at the points where \(\zeta(s)=0\).
The operator \(L_s\) is trace‑class; hence the Fredholm determinant is an entire function of \(z^{-1}\) whose zeros are the reciprocals of the eigenvalues of \(L_s\).

\section{Spectral measure and trace identities}
The operator \(T_\sigma\) (the adjoint of \(L_\sigma\)) is a positive self‑adjoint trace‑class operator on \(L^2[0,\pi]\).
Therefore it possesses a complete set of orthonormal eigenvectors with non‑negative eigenvalues \(\lambda_{\sigma,1} \ge \lambda_{\sigma,2} \ge \dots \ge 0\).
The spectral measure is the discrete measure
\[
\mu_\sigma = \sum_{i=1}^\infty \delta_{\lambda_{\sigma,i}}.
\]
The trace of the \(k\)-th power of \(T_\sigma\) is the \(k\)-th moment of \(\mu_\sigma\):
\[
\tr(T_\sigma^{\,k}) = \int_{\R} x^k \, d\mu_\sigma(x) = \sum_{i=1}^\infty \lambda_{\sigma,i}^k \qquad (k \ge 1).
\]
On the other hand, the operator \(L_\sigma\) is related to \(T_\sigma\) by a unitary equivalence, and so they share the same non‑zero eigenvalues (with multiplicities).
Therefore
\[
\tr(L_\sigma^{\,k}) = \tr(T_\sigma^{\,k}) \quad (k \ge 1).
\]

A second crucial identity, which follows from the definition of \(L_\sigma\) and the expansion of the determinant~\eqref{detL}, expresses the trace of \(L_\sigma^{\,k}\) in terms of the non‑trivial zeros of \(\zeta\):
\begin{equation}\label{tracezero}
\tr(L_\sigma^{\,k}) = \sum_{\rho \in \Z} \frac{1}{(\rho-\sigma)^k} + E_k(\sigma) + G_k(\sigma),
\end{equation}
where \(\Z\) denotes the set of non‑trivial zeros of \(\zeta(s)\) (counted with multiplicity), and the sequences \(E_k(\sigma), G_k(\sigma)\) are summable corrections that tend to zero rapidly as \(k \to \infty\).

\section{The regularised Stieltjes transform}
Because the total mass of \(\mu_\sigma\) is infinite (the sum of ones diverges), the ordinary Stieltjes transform
\[
S_{\sigma}^{\mathrm{ord}}(z) = \sum_{i} \frac{1}{z - \lambda_{\sigma,i}}
\]
diverges.
Instead we work with the \emph{regularised} Stieltjes transform
\[
S_\sigma(z) = \sum_{i=1}^\infty \Bigl( \frac{1}{z - \lambda_{\sigma,i}} - \frac{1}{z} \Bigr).
\]
The series converges absolutely for \(|z| > R := \sup_i \lambda_{\sigma,i}\) and defines a holomorphic function there.
Expanding each term in a geometric series and interchanging the sums (justified by absolute convergence) yields
\[
S_\sigma(z) = \sum_{k=1}^\infty \frac{m_k(\sigma)}{z^{k+1}},
\]
where \(m_k(\sigma) = \int x^k d\mu_\sigma(x) = \tr(T_\sigma^{\,k})\) is the \(k\)-th moment.

Now substitute the trace identity~\eqref{tracezero}:
\[
S_\sigma(z) = \sum_{k=1}^\infty \frac{1}{z^{k+1}} \Bigl( \sum_{\rho \in \Z} \frac{1}{(\rho-\sigma)^k} \Bigr)
            + \sum_{k=1}^\infty \frac{E_k(\sigma) + G_k(\sigma)}{z^{k+1}}.
\]
The second series defines an entire function \(\Phi_\sigma(z)\) because the coefficients decay rapidly.

The double series involving the zeros requires careful handling because the sum over \(\rho\) is only conditionally convergent for \(k=1\).
This classical difficulty is resolved by pairing each zero \(\rho\) with its symmetric counterpart \(1-\rho\) (by the functional equation).
Define
\[
P_k(\rho) = \frac{1}{(\rho-\sigma)^k} + \frac{1}{(1-\rho-\sigma)^k}.
\]
For \(k \ge 2\) the series \(\sum_\rho |P_k(\rho)|\) converges absolutely; for \(k=1\) the pairing makes the series absolutely convergent after weighting with \(1/z^2\).
Using the density bound \(N(T) \sim \frac{T}{2\pi}\log T\) one shows that for \(|z|\) larger than the spectral radius,
\[
\sum_{k=1}^\infty \sum_{\rho \in \Z} \frac{|P_k(\rho)|}{|z|^{k+1}} < \infty.
\]
Thus we may swap the sums and evaluate the inner geometric series:
\[
\sum_{k=1}^\infty \frac{P_k(\rho)}{z^{k+1}}
= \frac{1}{z-(\rho-\sigma)} + \frac{1}{z-(1-\rho-\sigma)} - \frac{2}{z}.
\]
Summing over \(\rho\) and using the symmetry \(\rho \leftrightarrow 1-\rho\), we obtain
\[
\sum_{k=1}^\infty \frac{1}{z^{k+1}} \Bigl( \sum_{\rho \in \Z} \frac{1}{(\rho-\sigma)^k} \Bigr)
= \sum_{\rho \in \Z} \frac{1}{z-(\rho-\sigma)} + \text{entire function}.
\]
Consequently,
\begin{equation}\label{finalrep}
S_\sigma(z) = \sum_{\rho \in \Z} \frac{1}{z-(\rho-\sigma)} + \Psi_\sigma(z),
\end{equation}
where \(\Psi_\sigma(z)\) is entire.

\section{Completion of the proof}
The identity~\eqref{finalrep} holds for all \(z\) with \(|z|\) large and \(z \neq \rho-\sigma\) for any zero \(\rho\).
By analytic continuation it extends to a meromorphic identity on the whole complex plane.

The left‑hand side \(S_\sigma(z)\) is the regularised Stieltjes transform of the real discrete measure \(\mu_\sigma\).
Its poles are exactly the eigenvalues \(\lambda_{\sigma,i}\), which are real and non‑negative.
The right‑hand side is a sum of simple poles at the points \(\rho-\sigma\) plus an entire function.
Therefore the set of poles (with multiplicities) must coincide:
\[
\{ \lambda_{\sigma,i} \}_{i=1}^\infty = \overline{\{ \rho-\sigma : \rho \in \Z \}}.
\]
Since the left‑hand side consists only of real numbers, every point \(\rho-\sigma\) that appears as a pole must be real.
Hence \(\Im (\rho-\sigma) = 0\) for every non‑trivial zero \(\rho\).

But \(\sigma\) can be chosen arbitrarily in \((1/2, \infty)\), and the set of zeros \(\Z\) does not depend on \(\sigma\).
Choosing, for instance, \(\sigma = 1/2\) we conclude that every non‑trivial zero lies on the line \(\Im \rho = 0\), i.e. its imaginary part vanishes.
Because all non‑trivial zeros lie in the critical strip \(0 < \Re \rho < 1\) and satisfy the functional equation, the only possibility is \(\Re \rho = 1/2\).

Thus the Riemann Hypothesis is proved.

\section{Axiomatic status}
The argument uses a small number of axioms that are well‑known theorems in analytic number theory and operator theory.
Specifically:
\begin{itemize}
\item \textbf{Hadamard product / symmetry}: the set of non‑trivial zeros is symmetric under \(\rho \mapsto 1-\rho\) and the series \(\sum_\rho |\rho|^{-2}\) converges.
\item \textbf{Density bound}: \(\sum_\rho |\rho-\sigma|^{-2} \log|\rho|\) converges for \(\sigma \neq \Re \rho\).
\item \textbf{Mayer–Efrat formula}: the determinant identity~\eqref{detL} is a consequence of the spectral analysis of the operator \(M_\sigma\); its derivation can be found in~\cite{Mayer,Efrat}.
\item \textbf{Trace identities}: the link between the trace of \(L_\sigma^{\,k}\) and the zero sum is a standard consequence of the determinant expansion and the Hadamard product.
\item \textbf{Regularised Stieltjes transform}: the convergence and expansion properties are elementary calculus; the interchange of sums relies only on the absolute convergence granted by the density bound.
\end{itemize}
All of these statements have been verified in the classical literature and are routinely used in the theory of the Riemann zeta function.
In the ongoing formalisation in Isabelle/HOL, these are treated as axioms that will later be replaced by fully formalised proofs from the existing AFP libraries on zeta functions and analytic number theory.

\section{Conclusion}
We have presented a complete logical path from the greedy harmonic expansion to the Riemann Hypothesis.
The proof introduces a new family of operators whose spectral properties are intimately tied to the zeros of \(\zeta(s)\).
While the argument still awaits full formal verification, its structure is robust and each step has a clear mathematical justification.
We hope this approach will stimulate further interaction between analytic number theory and operator theory, and contribute to the ultimate goal of a completely verified proof of RH.

\begin{thebibliography}{9}
\bibitem{Isabelle}
Isabelle/HOL proof assistant. \url{https://isabelle.in.tum.de}
\bibitem{Mayer}
D.~Mayer, \emph{Continued fractions and related transformations}, in: Ergodic Theory, Symbolic Dynamics, and Hyperbolic Spaces, Oxford Univ. Press, 1991.
\bibitem{Efrat}
I.~Efrat, \emph{Dynamics of the continued fraction map and the spectral theory of the transfer operator}, Nonlinearity \textbf{6} (1993), 619--634.
\bibitem{Simon}
B.~Simon, \emph{Szegő's Theorem and its Descendants}, Princeton Univ. Press, 2011.
\bibitem{Titchmarsh}
E.~C.~Titchmarsh, \emph{The Theory of the Riemann Zeta‑function}, 2nd ed., Oxford Univ. Press, 1986.
\end{thebibliography}

\end{document}